\begin{table*}[htbp!] 
\begin{center}
\caption{The ALPAKA sample. }
\label{tab:tab1}
\begin{tabular}{ccccccc}
\\
\hline\hline \noalign{\smallskip}
%\\
ID  &  Name & RA (deg) & Dec (deg) & redshift & Field/Survey & Notes \\
\noalign{\smallskip}
\hline
\noalign{\medskip}
1 & GOODS-S 15503 & 53.08205 & -27.83995 & 0.561 & GOODS-S & - \\ 
2 & COSMOS 2989680 & 150.43186 & 2.80261 & 0.623 & COSMOS & Minor merger\\
3 & COSMOS 1648673 & 149.98144 & 2.25321 & 1.445 & COSMOS & AGN(?) \\
4 & ALMA.03  & 333.99393 & -17.62988 & 1.453 & XCS & Cluster member \\
5 & ALMA.010 & 333.98853 & -17.63149 & 1.453 & XCS & Cluster member\\ %-co5field0
  %-co4field0 and co4field1
6 & ALMA.08 & 333.99266 & -17.63950 & 1.456 & XCS & Cluster member\\ % -co3field0
7 & ALMA.01 & 333.99233 & -17.63737 & 1.466 & XCS & Cluster member\\%
8 & ALMA.06 & 333.99880 & -17.63306 & 1.467 & XCS & Cluster member \\% -co2field1-co2field0
9 & ALMA.013 & 333.99909 & -17.63797 & 1.471 & XCS & Cluster member\\
10 & SHiZELS-19 & 149.79817 & 2.39008 & 1.484 & COSMOS & - \\
11 & SpARCS J0225-371 & 36.44191 & -3.92436 & 1.599 & SpARCS & Cluster member\\%
12 & SpARCS J0224-159 & 36.11320 & -3.40037 & 1.634 & SpARCS & Cluster member\\%
13 & COSMOS 3182 & 150.07594 & 2.21182 & 2.103 & COSMOS & Protocluster member\\
14 & Q2343-BX610 & 356.53934 & 12.82202 & 2.211 & SINS/zC-SINF & AGN(?)\\
15 & GS30274 & 53.13114 & -27.77319 & 2.225 & GOODS-S & AGN \\
16 & HXMM01-a & 35.06938 & -6.02830 & 2.311 & HerMES & Group member\\
17 & HXMM01-b+c & 35.06907 & -6.02904 & 2.308 & HerMES &  Group member, merger\\
18 & HATLAS J084933-W & 132.38994 & 2.24573 & 2.407 & H-ATLAS & AGN, protocluster member\\
19 & CLJ1001-131077 & 150.23728 & 2.33813 & 2.494 & COSMOS & Cluster member\\ %CO_1_3_3/CO_1_4_3 - SB1
20 & CLJ1001-130949 & 150.23691 & 2.33579 & 2.504 & COSMOS & Cluster member\\ %CO_1_3_2/CO_1_4_4 - MS1
21 & CLJ1001-130891 & 150.23986 & 2.33646 & 2.513 & COSMOS & Cluster member\\ %CO_1_4_1/CO_1_3_1 -SB2
22 & Gal3 & 150.33139 & 2.16239 & 2.935 & COSMOS & -\\
23 & ADF22.1 & 334.38507 & 0.29551 & 3.089 & SSA2 & AGN, protocluster member\\%CO_1
24 & ADF22.5 & 334.38416 & 0.29323 & 3.094 & SSA2 & Protocluster member \\%CO_3
25 & ADF22.7 & 334.38120 & 0.29946 & 3.088 & SSA2 & AGN, protocluster member \\%CO_2
26 & Gal5 & 149.87724 & 2.28388 & 3.339 & COSMOS & -\\
27 & Gal4  & 150.27827 & 2.25887 & 3.431 & COSMOS & -\\
28 & W0410-0913 & 62.54425 & -9.21812 & 3.630 & WISE & AGN, protocluster member\\
\noalign{\smallskip}
\hline
\end{tabular}
\tablefoot{Columns 1 and 2: ALPAKA ID and mostly used name for the galaxy. Columns 3 and 4: coordinates of the center used for the kinematic fitting in Sect.\ref{sec:kinematic}. Column 5: redshift from the systemic velocity obtained from the kinematic fitting. Column 6: field or survey where the galaxies lie (GOOD-S: \citet{Giavalisco_2004}, COSMOS: \citet{Scoville_2007}, XCS: \citet{Romer_2001}, SpARCS: \citet{Muzzin_2009}, SINS/zC-SINF: \citet{Forster_2018}, HerMES: \citet{Oliver_2012}, H-ATLAS: \citet{Eales_2010}, SSA2: \citet{Steidel_1998}, WISE: \citet{Wright_2010}). Column 7: indication about the environment and the presence of an AGN taken from the literature (see details in Sect. \ref{sec:details}). The presence of a "?" indicates that for the corresponding galaxy there are hints of the presence of an AGN ( Sect.~\ref{sec:details}).}
\end{center}
\end{table*}
